\section{Conclusiones}

El simulador cumple el flujo principal solicitado: los clientes registran sus
archivos en un servidor central, buscan por nombre, descargan archivos por
hash y tamano desde uno o varios pares, y pueden recurrir a una busqueda
distribuida cuando el servidor no entrega resultados. Las pruebas locales con
tres clientes demostraron busqueda centralizada, busqueda distribuida,
fallback automatico, descarga segmentada y tolerancia a la desaparicion de una
carpeta compartida.

La decision de usar TCP para todos los mensajes simplifico la confiabilidad del
transporte y permitio reutilizar el mismo framing. El costo es que cada
propagacion abre conexiones, por lo que el TTL debe mantenerse bajo. Para la
red de prueba, TTL 4 fue suficiente; subirlo a 6 solo tendria sentido en una
red mas grande y con mediciones de carga.

La principal limitacion es que el servidor central no elimina automaticamente
pares que ya no estan disponibles. Si un cliente se apaga sin avisar, sus
metadatos pueden permanecer hasta que falle una transferencia o una busqueda
distribuida encuentre alternativas. Como mejora futura se agregarian heartbeats
o expiracion de pares en el registro central, junto con una opcion explicita de
refrescar metadatos despues de cambios en la carpeta compartida.

Otra mejora seria enriquecer el protocolo distribuido con el puerto de datos
del remitente inmediato. La version actual conserva el protocolo congelado y
usa la cache de IDs vistos para evitar ciclos, lo cual funciona para la prueba,
pero un campo de remitente permitiria excluir con mas precision el enlace por
el que llego una consulta.
